\section{Diagnosis: what \vfour{} does against a strong opponent}
\label{sec:diagnosis}

This is the central section of the report. It establishes four things in order:
what the failure is, that it is not the failure the v0.4 study described, what
causes it, and what each contributing defect costs. Every number is produced by
an instrument described in \S\ref{sec:diag-instrument} and carries the artifact
it came from.

\subsection{The instrument}
\label{sec:diag-instrument}

The measurements below come from a diagnostic command added for this study,
\texttt{fish pathology}, together with five probes that re-derive its headline
quantities independently. Three definitions do the work.

\begin{definition}[Provably dead ask]
\label{def:dead}
An ask of card $c$ from target $t$ by actor $a$ is \emph{provably dead} if $a$'s
own public deduction state establishes that $t$ does not hold $c$: either $c$'s
holder is already known and is not $t$, or $t$ has been excluded from $c$'s
candidate set by a certificate. The predicate is a deduction from the actor's own
hand and the public record. No posterior and no policy prior enter it, so a
policy that filters on it cannot acquire a modelling error --- it can only lose
moves whose probability of success is exactly zero.
\end{definition}

\begin{definition}[Dead run]
A maximal sequence of consecutive asks in one game, every one of which is
provably dead in the sense of Definition~\ref{def:dead}.
\end{definition}

\begin{definition}[Starved turn]
A turn at which \emph{every} legal ask is provably dead, so the actor has no live
alternative. Starved turns separate a policy that chooses dead asks from a policy
that is forced into them.
\end{definition}

Two populations are used throughout: \emph{mirror} play, in which both teams run
the same policy, and the \emph{held-out} contrast against \vthree{}, which is the
population on which \vfour{} was evaluated for publication.

\subsection{The baseline, and why the published evaluation could not see it}
\label{sec:diag-baseline}

\begin{table}[t]
\centering
\caption{\vfour{} measured on the same instrument against itself and against
\vthree{}. Mirror: \numMirrorGames{} games, seed 31. Held-out: 600 games, seed
90210. Source: \texttt{research/v05/results/P0-v04-pathology.md}.}
\label{tab:baseline}
\begin{tabular}{lrr}
\toprule
& \textbf{mirror} & \textbf{vs.\ \vthree{}} \\
\midrule
events per game                       & \numMirrorEvents{}   & \numVthreeEvents{} \\
ask hit rate                          & \numMirrorHit\%      & \numVthreeHit\% \\
provably dead asks                    & \numDeadAsk\%        & \numVthreeDeadAsk\% \\
exact repeat (actor, card, target)    & \numRepeatAsk\%      & \numVthreeRepeat\% \\
longest dead run                      & \numDeadRunLongest{} & \numVthreeDeadRun{} \\
games with a dead run $\geq 6$        & \numDeadRunGames\%   & \numVthreeDeadRunGames\% \\
starved turns                         & \numStarved\%        & \numVthreeStarved\% \\
asks inside a half-suit the actor's team already owns & \numOwnLocked\% & \numVthreeOwnLocked\% \\
declarations wrong                    & \numDeclWrong\%      & \numVthreeDeclWrong\% \\
declarations at or after the forcing horizon & \numPostHorizonDecl{} & 0 \\
\quad of those, wrong                 & \numPostHorizonWrong\% & --- \\
forced-endgame declarations wrong     & \numForcedWrongRate\% & \numForcedWrongRate\% \\
\bottomrule
\end{tabular}
\end{table}

Table~\ref{tab:baseline} is the whole problem in one place. On the population
\vfour{} was published against, it is a clean policy: \numVthreeDeadAsk\% dead
asks, no game containing a dead run of six, no game reaching the forcing horizon
at all. Against a copy of itself, two in five of its asks name a card it can
prove the target does not hold, two in five are exact repetitions, and a third of
games contain an extended run of guaranteed misses. The two columns differ by a
factor of fourteen in dead-ask rate. A policy panel of weaker opponents cannot
detect this, and a competent human opponent is closer to the mirror column than
\vthree{} is --- which is how the failure was first reported, from live play, and
not from any experiment in the v0.4 study.

The pathology is not caused by the approximate inference path. Substituting the
exact reference belief reduces dead asks from \numDeadAsk\% to
\numBlockDeadAsk\% and leaves the longest dead run at \numBlockDeadRun{}. The
belief reports the hit probability correctly and the score prefers the move
anyway: this is a policy defect, not an inference defect.

\subsection{The deadlock is a two-question cycle}
\label{sec:diag-cycle}

Of \numScanGames{} mirror games scanned in the forensic probe, \numLongGames{}
exceed 300 events. In those games the mean longest dead run is \numCycleLen{}
asks, and the mean number of distinct (actor, card, target) triples inside that
run is \numCycleTriples{}: in \numPureCycle{} of \numLongGames{} the entire run
consists of exactly two questions, alternating, each asked well over a hundred
times. Two seats on opposite teams ask each other the same guaranteed-miss
question until an event counter stops them.

The reason a deterministic policy can do this is that the position is a fixed
point. Sampling each dumped run at its onset and at $+40$ and $+80$ events, every
observer's exact mean per-card certainty is identical to the last digit at all
three samples, as are the hand counts, the unresolved counts and the score. The
only quantity that changes is a counter: each repeated ask appends a fresh
ask-legality certificate that duplicates one already held, so the stored
certificate list grows without bound while the posterior it induces does not
move. Both policies re-derive the same argmax from the same information state,
and the cycle closes.

\subsection{The deadlock is not an information freeze}
\label{sec:diag-information}

The v0.4 termination study attributes non-termination to the locked-half-suit
theorem, stating that such a half-suit is frozen and that, for the same reason,
no further information about its allocation can ever arrive. That account fails
in two independent ways, and \S\ref{sec:corrections} retracts it.

\paragraph{First: the deadlock is mostly not in a locked half-suit.} A ground-truth
lock census taken at \emph{every} event of the dead run --- not only at its ends
--- finds no half-suit locked to any team at any point in the run in
\numNoLockPooled{} of \numNoLockPooledN{} long games pooled over four seeds
(\numNoLockPct\%). Measuring the stronger and more relevant condition, only
\numRunAsksLocked\% of dead-run asks sit inside a half-suit that is locked at that
moment, and only \numCycleInLock\% of long games have their entire cycle inside a
lock. In the majority of deadlocks the theorem has nothing to bite on: at the
onset of one traced run all nine half-suits are live and genuinely split,
\numTraceLive{} of \numTraceLegal{} legal asks are not provably dead, and the
policy ranks a guaranteed miss first and then repeats it \numTraceRepeats{}
times.

\paragraph{Second: where a half-suit \emph{is} locked, information still arrives.}
This is a statement about the rules before it is a measurement. An ask inside a
half-suit locked to the asker's own team is legal --- the set is active, the
target is a live opponent, the asker holds another card of the set and lacks the
one named --- and the deduction machinery publishes the asker's own exclusion
regardless of whether the ask succeeds. ``The asker does not hold this card'' is
public, permanent information about the allocation of a locked half-suit, and
with six cards split among three teammates it frequently resolves a card
outright.

Pooled over three independent seeds and \numCertN{} legal asks inside a
half-suit the asker's team in fact owns: \numCertRate\% strictly increase a
teammate's exact probability of the correct allocation, by a mean of
$+\numCertMean$ and a maximum of $+\numCertMax$. Re-measured with two statistics
that avoid the disputed estimator entirely --- the combinatorial support size of
the half-suit's cards, and the Shannon entropy of the exact marginals ---
\numEntropyRate\% strictly lower the entropy, by a mean of \numEntropyMean{}
nats. Asks \emph{outside} the locked half-suit shrink its support in
0 of \numOutsideSupport{} cases at any seed, which is the expected structural
result and confirms that the certificate lands where it is aimed. Playing such
asks greedily, a team's best allocation probability reaches certainty in
\numLadderRate\% of attempts in a median of \numLadderMedian{} asks.

Ownership monotonicity therefore does not imply informational monotonicity, and
the two must not be merged. The v0.4 \emph{paper} states this correctly as an
observation; the v0.4 termination artifact does not, and it is that artifact's
claim which \S\ref{sec:corrections} withdraws.

\paragraph{The leak, priced.} These certificates are public, so opponents read
them too. Over the same ladders, the team's allocation probability on the
half-suit it is trying to cash rises by $+\numTeamGain$ while the opponents' best
locked-half-suit allocation probability rises by $+\numOppGain$ --- an exchange
rate of about $\numLeakRatio\!:\!1$ in the sender's favour. The opponents' gain
is indirect, arriving through capacity coupling rather than through the locked
half-suit itself, which they can never take.

\paragraph{What the channel costs, and the qualification that matters.} A
guaranteed miss donates the turn. In \vfour{}'s own value function a donated turn
is worth \numDonationCost{} half-suits of final differential, and empirically the
receiving team's take from a miss inside a dead run is \numDonationDead{}
half-suits, with \numDonationNil\% of such donations yielding the opponents
nothing at all.
The qualification is important and was overstated in the first pass of this
analysis: the lock census that selects these asks is a ground-truth statistic. At
the deadlock states examined, the exact probability an owning team assigns to its
own ownership has mean \numOwnProbMean{} and maximum \numOwnProbMax{}, and not
one of the \numOwnProbPairs{} (observer, half-suit) pairs clears the
\numProvableThresh{} bar at which ownership would be provable. A \vfive{} ask rule therefore \emph{cannot} condition on ``a
half-suit my team provably owns''; the same conclusion is reached independently
in \S\ref{sec:diag-channels}, where the provable version of the condition holds
at \numProvableOwn\% of mirror decisions. The channel is real and it must be
priced against ownership probability, not gated on a certainty nobody has.

\subsection{The cause: ownership features that are not gated by hit probability}
\label{sec:diag-features}

The ask score is a twenty-feature linear function refined by a one-ply
expectimax. Four of its features measure ownership progress and none of them is
multiplied by the probability that the ask succeeds: own-set progress
($+\numOwnSetWeight$), team control ($+\numTeamControlWeight$), lock completion
($+\numLockWeight$) and the completion bonus ($+\numCompletionWeight$).
Evaluated at the compiled weights they pay up to $\numOwnershipPay$ at $p = 0$,
against $\numHitWeight \cdot p$ for the hit-probability term. A card the actor's whole team already owns
therefore outscores a genuine chance at a card it does not.

Three further features push the same way. The score's only explicit information
term is the binary entropy of the ask outcome, at weight $\numEntropyWeight$: the
ask whose result would resolve a full bit is penalised most, and the ask whose
result is already known is penalised not at all. This is a negative
value-of-information term in the literal sense --- the policy does not merely
ignore information, it charges for it. A second feature rewards re-asking in the
half-suit last asked in ($+\numRepeatWeight$), which is a direct mechanism for the
\numRepeatAsk\% exact-repeat rate. A third scales the exposure penalty by $(1-p)$
and is therefore maximised exactly when the ask is certain to miss.

Evaluating the compiled weight vector at two representative feature vectors, a
provably dead ask inside a half-suit the team already owns outright scores
\numDeadScore{} while a maximally informative fresh ask scores \numLiveScore{}.
That is an illustration rather than a game measurement --- the shipped policy
adds an expectimax term, so the argmax is not the linear argmax --- but the
decomposition of real decisions agrees with it. At the traced deadlock onsets the
gap between the dead ask ranked first and the best live ask decomposes into the
same three contributions every time: entropy, exposure-on-miss and target hand
size outvote the hit-probability term.

Two facts establish that this is a choice and not a constraint. Starved turns,
at which no live ask exists, are \numStarved\% of asks; \numVoluntaryDead\% of
provably dead asks are made with a live alternative available. And the most
informative legal ask available at a deadlock state ranks, in \vfour{}'s own
score, \numInfoRank{} out of roughly forty on average. At the \numInfoFreeStates{}
deadlock states examined the ask actually played moved no teammate's exact
posterior in every case, while an ask that would have moved it existed in every
case. The ask rule is not slightly mispriced on information; it does not contain
the term.

\paragraph{Causal isolation.} Filtering provably dead asks out of the candidate
set, with nothing else changed, collapses the longest dead run from
\numFilterBaseDeadRun{} to \numFilterDeadRun{}, takes games containing a dead run
of six from \numFilterBaseGames\% to 0\%, takes declarations at or after the
horizon from \numFilterBasePostHorizon{} to zero, takes misdeclarations from \numFilterBaseDeclWrong\% to
\numFilterDeclWrong\%, and scores \numFilterWin\% against the shipped policy ---
neutral. The single change removes the pathology at no measured cost in strength,
which identifies it as the load-bearing defect.

Removing the repetition without pricing information is a different matter, and it
fails. Forbidding exact repeats alone terminates every game, taking events per
game from \numNoRepeatBaseEvents{} to \numNoRepeatEvents{} and the longest dead
run from \numNoRepeatBaseRun{} to \numNoRepeatRun{}, but it loses to the shipped
policy at \numNoRepeatWin\% over \numNoRepeatGames{} games. Breaking the cycle is
not sufficient; the candidate set has to be restricted to asks that can actually
do something.

\subsection{Terminating by force, and terminating by patience}
\label{sec:diag-dilemma}

\vfour{} terminates deadlocks with a graduated escalation keyed on the public
event count. Raising the horizon so that the rule never fires is worth
\numPatientWin\% against the shipped configuration (deal-clustered
\numPatientWinCI\%, $n = \numPatientN$), with the patient side declaring at
\numPatientDeclAcc\% accuracy against the forcing side's \numForcingDeclAcc\%.
But pure patience does not terminate: the unforced mirror reaches the action cap
in \numPatientCap\% of games, with \numPatientDeadAsk\% of asks provably dead and
a longest run of \numPatientDeadRun{}. Table~\ref{tab:dilemma} states the
dichotomy.

\begin{table}[t]
\centering
\caption{The two settings \vfour{} has, and what each costs. Sources:
\texttt{research/v05/results/P0-v04-pathology.md},
\texttt{research/v05/results/P0b-forcing-dilemma.md}.}
\label{tab:dilemma}
\begin{tabular}{lrr}
\toprule
& \textbf{forcing (shipped)} & \textbf{patient} \\
\midrule
declarations wrong                        & \numDeclWrong\%     & \numPatientDeclWrong\% \\
games ending by action-cap adjudication   & 0\%                 & \numPatientCap\% \\
provably dead asks                        & \numDeadAsk\%       & \numPatientDeadAsk\% \\
longest dead run                          & \numDeadRunLongest{} & \numPatientDeadRun{} \\
\bottomrule
\end{tabular}
\end{table}

Note the third row. In the patient mirror more than half of all asks are already
provably dead: the turns are being burned either way. They are simply not being
spent on asks that emit useful certificates. That is the sense in which the
project owner's judgement --- that a risky ask is better than a guaranteed-zero
declaration --- is a statement about the ask policy and not about the
declaration rule.

\paragraph{The second stage of the escalation is an unconditional cash.} Past a
second horizon the rule returns ``declare'' before inspecting the allocation
probability at all, and simultaneously zeroes the team-ownership floor and the
marginal pre-gate. The consequence is measurable in the policy's own stated
confidence: bucketing declarations by the probability the policy itself attached
to them, \numLowConfDecl{} declarations in \numLowConfGames{} games are made at a
self-assessed probability below \numLowConfThresh{} and every one of them is
wrong. Pre-horizon declarations are wrong
\numPreHorizonWrong\%; at or after the horizon, \numHorizonWrong\%. Disabling the
second stage is worth \numStageTwoCost{} points against a mirror opponent and
changes nothing against \vthree{}, the detective, the lockout policy, the hunter
or the bluffer --- differences of at most a third of a point, in both directions.
It is a pure strong-opponent tax.

One hypothesis about this rule did \emph{not} hold and is recorded so that it is
not repeated: the horizon does not fire on healthy long games. Every one of the
\numHorizonGames{} games (\numHorizonReach\% of \numLowConfGames{}) that reached it was genuinely deadlocked,
and raising the deadlock threshold to a run of 80 leaves the split unchanged. The
cost is not when the rule fires but what it does when it fires.

\subsection{The forced endgame: an allocation with posterior probability zero}
\label{sec:diag-forced}

When one team is cardless the other must declare every remaining half-suit. Over
\numForcedGames{} distinct mirror games, all \numForcedN{} forced declarations
are wrong. An independently instrumented replay that shares no code with the
diagnostic probe and recomputes each declaration's correctness from the opening
deal reproduces the result at five further mirror seeds, disagreeing with the
engine's own flag on none of the declarations it re-scored.
This is not a hard decision problem and not an accounting artefact; it is one
mechanical defect with a deterministic consequence. The routine that names the
allocation picks each card's owner by an independent per-card argmax over
ownership marginals, with no capacity constraint. In every observed forced
endgame the truth splits the unresolved cards across two teammates and the argmax
piles them onto one. The named allocation therefore violates the declarer's own
capacity constraint, has exact posterior probability zero, and loses the
half-suit with certainty. All \numForcedN{} violate capacity; all
\numForcedN{} have exact posterior probability zero; the declarer's own stated
confidence is exactly zero in all \numForcedN{}; and the best \emph{feasible}
allocation would have been correct in \numFeasibleRight\% of them. The policy
misses by the minimum possible margin every time --- one card in
\numForcedMissRate\% of them, two in the remainder, never more.

The same defect disables the willingness ladder that selects the declarer. The
ladder thresholds a statistic that the unconstrained allocation drives to a hard
zero, and that statistic is exactly zero in \numRungZeroPairs{} of
\numRungZeroPairs{} surveyed (player, live half-suit) pairs. Four different
ladder shapes, including one with a bottom rung at \numLadderBottomRung{} and one with no rungs at
all, produce bit-identical output. The allocator is the defect; the ladder is
merely inert.

\subsection{The opponent model is two scalars, and one of them is decorative}
\label{sec:diag-prior}

\vfour{} reweights its rules posterior by a policy prior with two coefficients:
a weight on ``this seat asked in this half-suit'' and a weight on ``this seat took
turns but never asked here''. They are identical for all five other players and
never updated during a game, so a deliberately silent opponent is misread by
construction --- which is what the project owner exploited at the table.

Three measured findings qualify that story, and two of them contradict the
hypothesis this study began with.

\begin{enumerate}
  \item \textbf{Deleting the prior is worse, not better.} The policy-agnostic
        configuration loses to the shipped one by \numAgnosticCost{} points
        [\numAgnosticCI] in a paired ablation over a panel of three deceptive
        archetypes and the honest mirror, and by \numPriorDeleteCost{} points
        [\numPriorDeleteCI] against the three deceptive archetypes alone, which
        is the figure a claim about deception should quote. Per cell the effect
        is not uniform at the resolution these banks have: over
        \numPriorCells{} (seed, archetype) cells the ablation is worse in
        \numPriorCellsWorse{}, tied in \numPriorCellsTie{} and reversed in
        \numPriorCellsReversed{} (sign test $p \approx \numPriorSignP$).
        The exposure is over-weighting, not weighting: doubling the pair costs
        \numPriorDoubleCost{} points.
  \item \textbf{The silence channel does not exist.} Deleting the second
        coefficient entirely costs nothing measurable anywhere ---
        \numPhiWin\% [\numPhiCI] against the honest mirror, a paired delta of
        \numPhiPaired{} points over a four-opponent panel, and a fifteen-fold
        sweep that is flat inside its intervals. The algebra explains it: the
        exponent rearranges so that the second term is card-independent, and the
        capacity normalisation then removes it. The channel the owner's manoeuvre
        attacks is carrying no information to begin with.
  \item \textbf{The damage deception does is a signed, cell-local
        miscalibration.} Against a table of feinting opponents, a seat that has
        asked once in a half-suit really holds a given unresolved card of it
        \numFeintTrue\% of the time and the policy believes \numFeintBelief\% ---
        an over-confidence of $+\numFeintBias$, rising at two asks. Against
        honest opponents the same cells are calibrated to within \numFeintCalib{} and
        conservative. The exploit runs through the hard certificate, which the
        soft prior only amplifies; global belief accuracy is essentially
        untouched, mean absolute marginal error moving from \numBeliefErrHonest{}
        under an honest opponent to within $\pm\numBeliefErrTol$ of it under every archetype
        tested.
\end{enumerate}

No archetype converted its corruption of the policy's beliefs into wins: shipped
\vfour{} beats all three (\numVsFeint\%, \numVsSilent\%, \numVsWithholder\%).
The silent figure here is against the hit-probability-capped variant of that
archetype, which is the competent one and the one this subsection's belief
measurements use; \S\ref{sec:results-deception} runs the uncapped variant, on
which both policies score far higher, so the two figures are not the same
measurement. The
owner's live result is therefore not reproduced by commitment deception alone,
and \S\ref{sec:limitations} records what distinguishes the two settings.

\subsection{The value function carries one feature, and the mirror was never fitted against}
\label{sec:diag-value}

The sixteen-feature value model is, to measurement precision, one feature. Its
held-out game-clustered $R^2$ is \numValueRSq{}; the $R^2$ of a bias term plus
the score differential alone is \numValueRSqScore{}, which is higher. Two of the
sixteen features are algebraically the same feature, and five more cancel
identically at every call site --- moving all five to $\pm\numDeadFeatureSwing$ changes none of
\numDeadFeatureDecisions{} declaration decisions. Capacity is not the binding
constraint: a depth-3 gradient-boosted tree on the same rows reaches
\numValueRSqTree{}. And the model is least informative exactly where it is needed
most, reaching only \numValueRSqTail{} in the tail of games past \numTailEvents{}
events.

Splitting the model by use is the actionable result. Deleting the ask-side
expectimax entirely is worth \numAskSideValue{} [\numAskSideCI] --- nothing ---
while reverting declarations to fixed thresholds costs \numDeclSideValue{}
[\numDeclSideCI]. The part of the value model that earns its keep is the
declaration rule; the ask-side use, which is the side the deadlock lives on,
does not.

Finally, the fitting objective could not have caught any of this. The mirror was
absent from the fitting panel in every recorded round, and the recorded
per-opponent win rates never fall below \numPanelFloor{}, so no roughly even opponent was
ever present. The objective is a soft minimum at $\beta = \numSoftMinBeta$, which
over a panel whose win rates span about six points gives a maximum-to-minimum
gradient weight ratio of \numSoftMinRatio{}: it is a weighted mean wearing the
name of a minimum. That is the direct explanation for why a mirror pathology of
this size survived fitting.

\subsection{Channels the policy cannot express}
\label{sec:diag-channels}

Two capabilities are absent rather than mispriced.

\paragraph{The target dimension.} Over \numAskDecisions{} mirror ask decisions,
at \numTargetTwo\% of decisions at least two opponents are hard-indistinguishable
holders of the card actually asked for, carrying a mean of \numTargetBits{} free
bits per ask at a mean hit-probability spread of \numTargetSpread{}; against
\vthree{} the figures are \numTargetTwoVthree\% and \numTargetBitsVthree{} bits.
The dominant refinement term in the ask score discards the target's identity
outright, and every target-dependent quantity anywhere in the score is material
--- hit probability, the target's reply threat, hand size and exposure. Nothing
prices what the choice of target tells a teammate. A separate verification
qualifies the ``free'' in that description: the indistinguishability classes are
defined by hard consistency, and choosing within one is not free once the
capacity-marginal probabilities are compared, so the channel must be priced.
Turn routing is available on the same dimension --- at \numRouteTwo\% of mirror
decisions the actor could hand the turn to any of two or more chosen opponents by
a provably dead ask --- and the policy has one soft penalty for donating the turn
to a dangerous seat and no concept of choosing whom to donate it to.

\paragraph{A knowledge model of anybody else.} The policy builds deduction state
only for itself; the only such objects it ever constructs are clones of its own.
Every technique that turns on what another seat knows --- teammate-directed
signalling, blackballing, choosing the best-informed declarer, baiting --- is out
of reach not because a feature is missing but because the state it would read
does not exist. Since the transcript is public, this is nearly free to add:
seat $j$'s knowledge is the public deduction state plus a posterior over $j$'s
hand.

\subsection{Summary of confirmed defects}
\label{sec:diag-summary}

\begin{table}[t]
\centering
\caption{Confirmed defects, each with the measurement that establishes it and the
mechanism of \S\ref{sec:mechanisms} that addresses it. Rows marked
``\emph{designed}'' are specified but not built.}
\label{tab:defects}
\begin{tabular}{p{0.30\textwidth}p{0.42\textwidth}l}
\toprule
\textbf{Defect} & \textbf{Measured cost} & \textbf{Fix} \\
\midrule
Ownership features not gated by hit probability &
\numDeadAsk\% dead asks, \numVoluntaryDead\% of them voluntary; longest dead run
\numDeadRunLongest{} & M1 \\
Unconstrained per-card allocation guess &
\numForcedWrongRate\% of \numForcedN{} distinct forced declarations wrong; willingness
ladder inert at \numRungZeroPairs{} of \numRungZeroPairs{} pairs & M2 \\
Escalation stage 2 cashes without inspecting probability &
\numStageTwoCost{} points against a mirror opponent, zero against every weaker
style & M8 \\
Negative value-of-information term &
weight $\numEntropyWeight$ on outcome entropy; most informative ask ranks
\numInfoRank{} of $\approx$40 & M3 (\emph{designed}) \\
Refinement term discards the target's identity &
\numTargetBits{} free bits per ask at \numTargetTwo\% of decisions & M5
(\emph{designed}) \\
No knowledge model of other seats &
every ``what does he know'' technique unreachable & M4 (\emph{designed}) \\
Deterministic argmax ask selection &
leakage pinned at the deterministic bound \cite{alvim-leakage};
\numRepeatAsk\% exact repeats & M6 (\emph{designed}) \\
Opponent model is two fixed global scalars &
silence channel measurably inert; feint miscalibration $+\numFeintBias$ & M7
(\emph{designed}) \\
Value model is effectively one feature &
held-out $R^2$ \numValueRSq{} against \numValueRSqScore{} for bias plus score
differential & M9 (\emph{designed}) \\
Mirror absent from the fitting panel &
soft-minimum gradient ratio \numSoftMinRatio{} at $\beta=\numSoftMinBeta$ & M10
(\emph{designed}) \\
\bottomrule
\end{tabular}
\end{table}
